
\documentclass{article}
\usepackage{colortbl}
\usepackage{makecell}
\usepackage{multirow}
\usepackage{supertabular}

\begin{document}

\newcounter{utterance}

\centering \large Interaction Transcript for game `matchit\_ascii', experiment `similar\_grid\_2', episode 6 with Qwen3.6{-}35B{-}A3B{-}FP8{-}without{-}reasoning.
\vspace{24pt}

{ \footnotesize  \setcounter{utterance}{1}
\setlength{\tabcolsep}{0pt}
\begin{supertabular}{c@{$\;$}|p{.15\linewidth}@{}p{.15\linewidth}p{.15\linewidth}p{.15\linewidth}p{.15\linewidth}p{.15\linewidth}}
    \# & $\;$A & \multicolumn{4}{c}{Game Master} & $\;\:$B\\
    \hline

    \theutterance \stepcounter{utterance}  
    & & \multicolumn{4}{p{0.6\linewidth}}{
        \cellcolor[rgb]{0.9,0.9,0.9}{
            \makecell[{{p{\linewidth}}}]{
                \texttt{\tiny{[P1$\langle$GM]}}
                \texttt{Συμμετέχετε σε ένα συνεργατικό παιχνίδι μαντεψιάς.} \\
\texttt{Ο στόχος είναι να βρείτε αν αυτό το πλέγμα και μια άλλη εικόνα που μόνο εγώ μπορώ να δω, είναι τα ίδια.} \\
\texttt{Παρακαλώ περιγράψτε πρώτα το πλέγμα σας. Στη συνέχεια, θα δώσω τη δική μου περιγραφή και μπορούμε να κάνουμε ο ένας στον άλλο ερωτήσεις σχετικά με τα πλέγματα για να καταλάβουμε αν είναι τα ίδια. Ξεκινήστε τώρα τη σύντομη περιγραφή του πλέγματός σας με "ΠΕΡΙΓΡΑΦΗ:" ακολουθούμενη από την περιγραφή. Μην προσθέσετε τίποτα άλλο. Αυτό είναι το πλέγμα σας:} \\
\texttt{X ▢ ▢ ▢ X} \\
\texttt{▢ X ▢ X ▢} \\
\texttt{▢ ▢ X ▢ ▢} \\
\texttt{▢ ▢ X ▢ ▢} \\
\texttt{▢ ▢ X ▢ ▢} \\
            }
        }
    }
    & \\ \\

    \theutterance \stepcounter{utterance}  
    & \multicolumn{4}{p{0.6\linewidth}}{
        \cellcolor[rgb]{0.9,0.9,0.9}{
            \makecell[{{p{\linewidth}}}]{
                \texttt{\tiny{[P1$\rangle$GM]}}
                \texttt{ΠΕΡΙΓΡΑΦΗ: Το πλέγμα είναι 5x5. Οι θέσεις με 'X' σχηματίζουν ένα σχήμα που μοιάζει με βέλος ή βελόνα που δείχνει προς τα πάνω. Συγκεκριμένα, υπάρχει μια κάθετη γραμμή από τρία 'X' στη μεσαία στήλη (θέσεις 3,3· 4,3· 5,3) και δύο επιπλέον 'X' στις γωνίες της πρώτης γραμμής (θέσεις 1,1 και 1,5).} \\
            }
        }
    }
    & & \\ \\

    \theutterance \stepcounter{utterance}  
    & & & \multicolumn{2}{p{0.3\linewidth}}{
        \cellcolor[rgb]{0.9,0.9,0.9}{
            \makecell[{{p{\linewidth}}}]{
                \texttt{\tiny{[GM$|$GM]}}
                \texttt{continue} \\
            }
        }
    }
    & & \\ \\

    \theutterance \stepcounter{utterance}  
    & & \multicolumn{4}{p{0.6\linewidth}}{
        \cellcolor[rgb]{0.9,0.9,0.9}{
            \makecell[{{p{\linewidth}}}]{
                \texttt{\tiny{[GM$\rangle$P2]}}
                \texttt{Συμμετέχετε σε ένα συνεργατικό παιχνίδι μαντεψιάς.} \\
\texttt{Ο στόχος είναι να βρείτε αν αυτό το πλέγμα και μια άλλη εικόνα που μόνο εγώ μπορώ να δω, είναι τα ίδια.} \\
\texttt{Παρακαλώ περιγράψτε πρώτα το πλέγμα σας. Στη συνέχεια, θα δώσω τη δική μου περιγραφή και μπορούμε να κάνουμε ο ένας στον άλλο ερωτήσεις σχετικά με τα πλέγματα για να καταλάβουμε αν είναι τα ίδια. Ξεκινήστε τώρα τη σύντομη περιγραφή του πλέγματός σας με "ΠΕΡΙΓΡΑΦΗ:" ακολουθούμενη από την περιγραφή. Μην προσθέσετε τίποτα άλλο. Αυτό είναι το πλέγμα σας:} \\
\texttt{X ▢ ▢ ▢ X} \\
\texttt{▢ ▢ ▢ X ▢} \\
\texttt{▢ ▢ X ▢ ▢} \\
\texttt{▢ ▢ X ▢ ▢} \\
\texttt{▢ ▢ ▢ ▢ ▢} \\
            }
        }
    }
    & \\ \\

    \theutterance \stepcounter{utterance}  
    & & & \multicolumn{4}{p{0.6\linewidth}}{
        \cellcolor[rgb]{0.9,0.9,0.9}{
            \makecell[{{p{\linewidth}}}]{
                \texttt{\tiny{[GM$\langle$P2]}}
                \texttt{ΠΕΡΙΓΡΑΦΗ: Το πλέγμα είναι 5x5. Υπάρχουν δύο 'X' στην πρώτη γραμμή, στις θέσεις (1,1) και (1,5). Υπάρχει ένα 'X' στη δεύτερη γραμμή, στη θέση (2,4). Υπάρχουν δύο 'X' στην τρίτη γραμμή, στη θέση (3,3), και δύο 'X' στην τέταρτη γραμμή, στη θέση (4,3). Η πέμπτη γραμμή είναι κενή.} \\
            }
        }
    }
     \\ \\

    \theutterance \stepcounter{utterance}  
    & & & \multicolumn{2}{p{0.3\linewidth}}{
        \cellcolor[rgb]{0.9,0.9,0.9}{
            \makecell[{{p{\linewidth}}}]{
                \texttt{\tiny{[GM$|$GM]}}
                \texttt{continue} \\
            }
        }
    }
    & & \\ \\

    \theutterance \stepcounter{utterance}  
    & & \multicolumn{4}{p{0.6\linewidth}}{
        \cellcolor[rgb]{0.9,0.9,0.9}{
            \makecell[{{p{\linewidth}}}]{
                \texttt{\tiny{[GM$\rangle$P2]}}
                \texttt{Αυτή είναι η ΠΕΡΙΓΡΑΦΗ μουΠΕΡΙΓΡΑΦΗ: Το πλέγμα είναι 5x5. Οι θέσεις με 'X' σχηματίζουν ένα σχήμα που μοιάζει με βέλος ή βελόνα που δείχνει προς τα πάνω. Συγκεκριμένα, υπάρχει μια κάθετη γραμμή από τρία 'X' στη μεσαία στήλη (θέσεις 3,3· 4,3· 5,3) και δύο επιπλέον 'X' στις γωνίες της πρώτης γραμμής (θέσεις 1,1 και 1,5).} \\
\texttt{Τώρα κάντε μια ερώτηση για να μάθετε νέες πτυχές του πλέγματός μου που μπορεί να διαφέρουν από το δικό σας πλέγμα. Ξεκινήστε με "ΕΡΩΤΗΣΗ:" και μην προσθέσετε τίποτα άλλο.} \\
            }
        }
    }
    & \\ \\

    \theutterance \stepcounter{utterance}  
    & & & \multicolumn{4}{p{0.6\linewidth}}{
        \cellcolor[rgb]{0.9,0.9,0.9}{
            \makecell[{{p{\linewidth}}}]{
                \texttt{\tiny{[GM$\langle$P2]}}
                \texttt{ΕΡΩΤΗΣΗ: Στο πλέγμα σας, υπάρχουν 'X' στη θέση (2,4) ή στη θέση (5,3);} \\
            }
        }
    }
     \\ \\

    \theutterance \stepcounter{utterance}  
    & & & \multicolumn{2}{p{0.3\linewidth}}{
        \cellcolor[rgb]{0.9,0.9,0.9}{
            \makecell[{{p{\linewidth}}}]{
                \texttt{\tiny{[GM$|$GM]}}
                \texttt{continue} \\
            }
        }
    }
    & & \\ \\

    \theutterance \stepcounter{utterance}  
    & & \multicolumn{4}{p{0.6\linewidth}}{
        \cellcolor[rgb]{0.9,0.9,0.9}{
            \makecell[{{p{\linewidth}}}]{
                \texttt{\tiny{[P1$\langle$GM]}}
                \texttt{Αυτή είναι η ΠΕΡΙΓΡΑΦΗ μουΠΕΡΙΓΡΑΦΗ: Το πλέγμα είναι 5x5. Υπάρχουν δύο 'X' στην πρώτη γραμμή, στις θέσεις (1,1) και (1,5). Υπάρχει ένα 'X' στη δεύτερη γραμμή, στη θέση (2,4). Υπάρχουν δύο 'X' στην τρίτη γραμμή, στη θέση (3,3), και δύο 'X' στην τέταρτη γραμμή, στη θέση (4,3). Η πέμπτη γραμμή είναι κενή.} \\
\texttt{ΕΡΩΤΗΣΗ: Στο πλέγμα σας, υπάρχουν 'X' στη θέση (2,4) ή στη θέση (5,3);Ξεκινήστε την απάντησή σας με "ΑΠΑΝΤΗΣΗ:" και μην προσθέσετε τίποτα άλλο.} \\
            }
        }
    }
    & \\ \\

    \theutterance \stepcounter{utterance}  
    & \multicolumn{4}{p{0.6\linewidth}}{
        \cellcolor[rgb]{0.9,0.9,0.9}{
            \makecell[{{p{\linewidth}}}]{
                \texttt{\tiny{[P1$\rangle$GM]}}
                \texttt{ΑΠΑΝΤΗΣΗ: Όχι, δεν υπάρχουν 'X' στη θέση (2,4) ούτε στη θέση (5,3). Στη θέση (2,4) υπάρχει κενό (▢) και στη θέση (5,3) υπάρχει 'X'.} \\
            }
        }
    }
    & & \\ \\

    \theutterance \stepcounter{utterance}  
    & & & \multicolumn{2}{p{0.3\linewidth}}{
        \cellcolor[rgb]{0.9,0.9,0.9}{
            \makecell[{{p{\linewidth}}}]{
                \texttt{\tiny{[GM$|$GM]}}
                \texttt{continue} \\
            }
        }
    }
    & & \\ \\

    \theutterance \stepcounter{utterance}  
    & & \multicolumn{4}{p{0.6\linewidth}}{
        \cellcolor[rgb]{0.9,0.9,0.9}{
            \makecell[{{p{\linewidth}}}]{
                \texttt{\tiny{[P1$\langle$GM]}}
                \texttt{Τώρα κάντε μια ερώτηση για να μάθετε νέες πτυχές του πλέγματός μου που μπορεί να διαφέρουν από το δικό σας πλέγμα. Ξεκινήστε με "ΕΡΩΤΗΣΗ:" και μην προσθέσετε τίποτα άλλο.} \\
            }
        }
    }
    & \\ \\

    \theutterance \stepcounter{utterance}  
    & \multicolumn{4}{p{0.6\linewidth}}{
        \cellcolor[rgb]{0.9,0.9,0.9}{
            \makecell[{{p{\linewidth}}}]{
                \texttt{\tiny{[P1$\rangle$GM]}}
                \texttt{ΕΡΩΤΗΣΗ: Στο πλέγμα σας, υπάρχει 'X' στη θέση (3,3);} \\
            }
        }
    }
    & & \\ \\

    \theutterance \stepcounter{utterance}  
    & & & \multicolumn{2}{p{0.3\linewidth}}{
        \cellcolor[rgb]{0.9,0.9,0.9}{
            \makecell[{{p{\linewidth}}}]{
                \texttt{\tiny{[GM$|$GM]}}
                \texttt{continue} \\
            }
        }
    }
    & & \\ \\

    \theutterance \stepcounter{utterance}  
    & & \multicolumn{4}{p{0.6\linewidth}}{
        \cellcolor[rgb]{0.9,0.9,0.9}{
            \makecell[{{p{\linewidth}}}]{
                \texttt{\tiny{[GM$\rangle$P2]}}
                \texttt{ΑΠΑΝΤΗΣΗ: Όχι, δεν υπάρχουν 'X' στη θέση (2,4) ούτε στη θέση (5,3). Στη θέση (2,4) υπάρχει κενό (▢) και στη θέση (5,3) υπάρχει 'X'.} \\
\texttt{ΕΡΩΤΗΣΗ: Στο πλέγμα σας, υπάρχει 'X' στη θέση (3,3);Ξεκινήστε την απάντησή σας με "ΑΠΑΝΤΗΣΗ:" και μην προσθέσετε τίποτα άλλο.} \\
            }
        }
    }
    & \\ \\

    \theutterance \stepcounter{utterance}  
    & & & \multicolumn{4}{p{0.6\linewidth}}{
        \cellcolor[rgb]{0.9,0.9,0.9}{
            \makecell[{{p{\linewidth}}}]{
                \texttt{\tiny{[GM$\langle$P2]}}
                \texttt{ΑΠΑΝΤΗΣΗ: Ναι, υπάρχει 'X' στη θέση (3,3).} \\
            }
        }
    }
     \\ \\

    \theutterance \stepcounter{utterance}  
    & & & \multicolumn{2}{p{0.3\linewidth}}{
        \cellcolor[rgb]{0.9,0.9,0.9}{
            \makecell[{{p{\linewidth}}}]{
                \texttt{\tiny{[GM$|$GM]}}
                \texttt{continue} \\
            }
        }
    }
    & & \\ \\

    \theutterance \stepcounter{utterance}  
    & & \multicolumn{4}{p{0.6\linewidth}}{
        \cellcolor[rgb]{0.9,0.9,0.9}{
            \makecell[{{p{\linewidth}}}]{
                \texttt{\tiny{[GM$\rangle$P2]}}
                \texttt{Τώρα κάντε μια ερώτηση για να μάθετε νέες πτυχές του πλέγματός μου που μπορεί να διαφέρουν από το δικό σας πλέγμα. Ξεκινήστε με "ΕΡΩΤΗΣΗ:" και μην προσθέσετε τίποτα άλλο.} \\
            }
        }
    }
    & \\ \\

    \theutterance \stepcounter{utterance}  
    & & & \multicolumn{4}{p{0.6\linewidth}}{
        \cellcolor[rgb]{0.9,0.9,0.9}{
            \makecell[{{p{\linewidth}}}]{
                \texttt{\tiny{[GM$\langle$P2]}}
                \texttt{ΕΡΩΤΗΣΗ: Στο πλέγμα σας, υπάρχει 'X' στη θέση (2,4);} \\
            }
        }
    }
     \\ \\

    \theutterance \stepcounter{utterance}  
    & & & \multicolumn{2}{p{0.3\linewidth}}{
        \cellcolor[rgb]{0.9,0.9,0.9}{
            \makecell[{{p{\linewidth}}}]{
                \texttt{\tiny{[GM$|$GM]}}
                \texttt{continue} \\
            }
        }
    }
    & & \\ \\

    \theutterance \stepcounter{utterance}  
    & & \multicolumn{4}{p{0.6\linewidth}}{
        \cellcolor[rgb]{0.9,0.9,0.9}{
            \makecell[{{p{\linewidth}}}]{
                \texttt{\tiny{[P1$\langle$GM]}}
                \texttt{ΑΠΑΝΤΗΣΗ: Ναι, υπάρχει 'X' στη θέση (3,3).} \\
\texttt{ΕΡΩΤΗΣΗ: Στο πλέγμα σας, υπάρχει 'X' στη θέση (2,4);Ξεκινήστε την απάντησή σας με "ΑΠΑΝΤΗΣΗ:" και μην προσθέσετε τίποτα άλλο.} \\
            }
        }
    }
    & \\ \\

    \theutterance \stepcounter{utterance}  
    & \multicolumn{4}{p{0.6\linewidth}}{
        \cellcolor[rgb]{0.9,0.9,0.9}{
            \makecell[{{p{\linewidth}}}]{
                \texttt{\tiny{[P1$\rangle$GM]}}
                \texttt{ΑΠΑΝΤΗΣΗ: Όχι, δεν υπάρχει 'X' στη θέση (2,4).} \\
            }
        }
    }
    & & \\ \\

    \theutterance \stepcounter{utterance}  
    & & & \multicolumn{2}{p{0.3\linewidth}}{
        \cellcolor[rgb]{0.9,0.9,0.9}{
            \makecell[{{p{\linewidth}}}]{
                \texttt{\tiny{[GM$|$GM]}}
                \texttt{continue} \\
            }
        }
    }
    & & \\ \\

    \theutterance \stepcounter{utterance}  
    & & \multicolumn{4}{p{0.6\linewidth}}{
        \cellcolor[rgb]{0.9,0.9,0.9}{
            \makecell[{{p{\linewidth}}}]{
                \texttt{\tiny{[P1$\langle$GM]}}
                \texttt{Τώρα κάντε μια ερώτηση για να μάθετε νέες πτυχές του πλέγματός μου που μπορεί να διαφέρουν από το δικό σας πλέγμα. Ξεκινήστε με "ΕΡΩΤΗΣΗ:" και μην προσθέσετε τίποτα άλλο.} \\
            }
        }
    }
    & \\ \\

    \theutterance \stepcounter{utterance}  
    & \multicolumn{4}{p{0.6\linewidth}}{
        \cellcolor[rgb]{0.9,0.9,0.9}{
            \makecell[{{p{\linewidth}}}]{
                \texttt{\tiny{[P1$\rangle$GM]}}
                \texttt{ΕΡΩΤΗΣΗ: Στο πλέγμα σας, υπάρχει 'X' στη θέση (2,2);} \\
            }
        }
    }
    & & \\ \\

    \theutterance \stepcounter{utterance}  
    & & & \multicolumn{2}{p{0.3\linewidth}}{
        \cellcolor[rgb]{0.9,0.9,0.9}{
            \makecell[{{p{\linewidth}}}]{
                \texttt{\tiny{[GM$|$GM]}}
                \texttt{continue} \\
            }
        }
    }
    & & \\ \\

    \theutterance \stepcounter{utterance}  
    & & \multicolumn{4}{p{0.6\linewidth}}{
        \cellcolor[rgb]{0.9,0.9,0.9}{
            \makecell[{{p{\linewidth}}}]{
                \texttt{\tiny{[GM$\rangle$P2]}}
                \texttt{ΑΠΑΝΤΗΣΗ: Όχι, δεν υπάρχει 'X' στη θέση (2,4).} \\
\texttt{ΕΡΩΤΗΣΗ: Στο πλέγμα σας, υπάρχει 'X' στη θέση (2,2);Ξεκινήστε την απάντησή σας με "ΑΠΑΝΤΗΣΗ:" και μην προσθέσετε τίποτα άλλο.} \\
            }
        }
    }
    & \\ \\

    \theutterance \stepcounter{utterance}  
    & & & \multicolumn{4}{p{0.6\linewidth}}{
        \cellcolor[rgb]{0.9,0.9,0.9}{
            \makecell[{{p{\linewidth}}}]{
                \texttt{\tiny{[GM$\langle$P2]}}
                \texttt{ΑΠΑΝΤΗΣΗ: Ναι, υπάρχει 'X' στη θέση (2,2).} \\
            }
        }
    }
     \\ \\

    \theutterance \stepcounter{utterance}  
    & & & \multicolumn{2}{p{0.3\linewidth}}{
        \cellcolor[rgb]{0.9,0.9,0.9}{
            \makecell[{{p{\linewidth}}}]{
                \texttt{\tiny{[GM$|$GM]}}
                \texttt{continue} \\
            }
        }
    }
    & & \\ \\

    \theutterance \stepcounter{utterance}  
    & & \multicolumn{4}{p{0.6\linewidth}}{
        \cellcolor[rgb]{0.9,0.9,0.9}{
            \makecell[{{p{\linewidth}}}]{
                \texttt{\tiny{[GM$\rangle$P2]}}
                \texttt{Τώρα κάντε μια ερώτηση για να μάθετε νέες πτυχές του πλέγματός μου που μπορεί να διαφέρουν από το δικό σας πλέγμα. Ξεκινήστε με "ΕΡΩΤΗΣΗ:" και μην προσθέσετε τίποτα άλλο.} \\
            }
        }
    }
    & \\ \\

    \theutterance \stepcounter{utterance}  
    & & & \multicolumn{4}{p{0.6\linewidth}}{
        \cellcolor[rgb]{0.9,0.9,0.9}{
            \makecell[{{p{\linewidth}}}]{
                \texttt{\tiny{[GM$\langle$P2]}}
                \texttt{ΕΡΩΤΗΣΗ: Στο πλέγμα σας, υπάρχει 'X' στη θέση (2,3);} \\
            }
        }
    }
     \\ \\

    \theutterance \stepcounter{utterance}  
    & & & \multicolumn{2}{p{0.3\linewidth}}{
        \cellcolor[rgb]{0.9,0.9,0.9}{
            \makecell[{{p{\linewidth}}}]{
                \texttt{\tiny{[GM$|$GM]}}
                \texttt{continue} \\
            }
        }
    }
    & & \\ \\

    \theutterance \stepcounter{utterance}  
    & & \multicolumn{4}{p{0.6\linewidth}}{
        \cellcolor[rgb]{0.9,0.9,0.9}{
            \makecell[{{p{\linewidth}}}]{
                \texttt{\tiny{[P1$\langle$GM]}}
                \texttt{ΑΠΑΝΤΗΣΗ: Ναι, υπάρχει 'X' στη θέση (2,2).} \\
\texttt{ΕΡΩΤΗΣΗ: Στο πλέγμα σας, υπάρχει 'X' στη θέση (2,3);Ξεκινήστε την απάντησή σας με "ΑΠΑΝΤΗΣΗ:" και μην προσθέσετε τίποτα άλλο.} \\
            }
        }
    }
    & \\ \\

    \theutterance \stepcounter{utterance}  
    & \multicolumn{4}{p{0.6\linewidth}}{
        \cellcolor[rgb]{0.9,0.9,0.9}{
            \makecell[{{p{\linewidth}}}]{
                \texttt{\tiny{[P1$\rangle$GM]}}
                \texttt{ΑΠΑΝΤΗΣΗ: Όχι, δεν υπάρχει 'X' στη θέση (2,3).} \\
            }
        }
    }
    & & \\ \\

    \theutterance \stepcounter{utterance}  
    & & & \multicolumn{2}{p{0.3\linewidth}}{
        \cellcolor[rgb]{0.9,0.9,0.9}{
            \makecell[{{p{\linewidth}}}]{
                \texttt{\tiny{[GM$|$GM]}}
                \texttt{continue} \\
            }
        }
    }
    & & \\ \\

    \theutterance \stepcounter{utterance}  
    & & \multicolumn{4}{p{0.6\linewidth}}{
        \cellcolor[rgb]{0.9,0.9,0.9}{
            \makecell[{{p{\linewidth}}}]{
                \texttt{\tiny{[P1$\langle$GM]}}
                \texttt{Τώρα κάντε μια ερώτηση για να μάθετε νέες πτυχές του πλέγματός μου που μπορεί να διαφέρουν από το δικό σας πλέγμα. Ξεκινήστε με "ΕΡΩΤΗΣΗ:" και μην προσθέσετε τίποτα άλλο.} \\
            }
        }
    }
    & \\ \\

    \theutterance \stepcounter{utterance}  
    & \multicolumn{4}{p{0.6\linewidth}}{
        \cellcolor[rgb]{0.9,0.9,0.9}{
            \makecell[{{p{\linewidth}}}]{
                \texttt{\tiny{[P1$\rangle$GM]}}
                \texttt{ΕΡΩΤΗΣΗ: Στο πλέγμα σας, υπάρχει 'X' στη θέση (2,1);} \\
            }
        }
    }
    & & \\ \\

    \theutterance \stepcounter{utterance}  
    & & & \multicolumn{2}{p{0.3\linewidth}}{
        \cellcolor[rgb]{0.9,0.9,0.9}{
            \makecell[{{p{\linewidth}}}]{
                \texttt{\tiny{[GM$|$GM]}}
                \texttt{continue} \\
            }
        }
    }
    & & \\ \\

    \theutterance \stepcounter{utterance}  
    & & \multicolumn{4}{p{0.6\linewidth}}{
        \cellcolor[rgb]{0.9,0.9,0.9}{
            \makecell[{{p{\linewidth}}}]{
                \texttt{\tiny{[GM$\rangle$P2]}}
                \texttt{ΑΠΑΝΤΗΣΗ: Όχι, δεν υπάρχει 'X' στη θέση (2,3).} \\
\texttt{ΕΡΩΤΗΣΗ: Στο πλέγμα σας, υπάρχει 'X' στη θέση (2,1);Ξεκινήστε την απάντησή σας με "ΑΠΑΝΤΗΣΗ:" και μην προσθέσετε τίποτα άλλο.} \\
            }
        }
    }
    & \\ \\

    \theutterance \stepcounter{utterance}  
    & & & \multicolumn{4}{p{0.6\linewidth}}{
        \cellcolor[rgb]{0.9,0.9,0.9}{
            \makecell[{{p{\linewidth}}}]{
                \texttt{\tiny{[GM$\langle$P2]}}
                \texttt{ΑΠΑΝΤΗΣΗ: Όχι, δεν υπάρχει 'X' στη θέση (2,1).} \\
            }
        }
    }
     \\ \\

    \theutterance \stepcounter{utterance}  
    & & & \multicolumn{2}{p{0.3\linewidth}}{
        \cellcolor[rgb]{0.9,0.9,0.9}{
            \makecell[{{p{\linewidth}}}]{
                \texttt{\tiny{[GM$|$GM]}}
                \texttt{continue} \\
            }
        }
    }
    & & \\ \\

    \theutterance \stepcounter{utterance}  
    & & \multicolumn{4}{p{0.6\linewidth}}{
        \cellcolor[rgb]{0.9,0.9,0.9}{
            \makecell[{{p{\linewidth}}}]{
                \texttt{\tiny{[GM$\rangle$P2]}}
                \texttt{Τώρα πάρε μια απόφαση. Τι πιστεύεις: είναι το δικό σου πλέγμα και το άλλο πλέγμα που περιγράφεται το ίδιο; Γράψε "ΑΠΟΦΑΣΗ: ίδιο πλέγμα" αν πιστεύεις ότι είναι το ίδιο πλέγμα ή "ΑΠΟΦΑΣΗ: διαφορετικά πλέγματα" αν πιστεύεις ότι είναι διαφορετικά πλέγματα. Μην προσθέσεις τίποτα άλλο.} \\
            }
        }
    }
    & \\ \\

    \theutterance \stepcounter{utterance}  
    & & & \multicolumn{4}{p{0.6\linewidth}}{
        \cellcolor[rgb]{0.9,0.9,0.9}{
            \makecell[{{p{\linewidth}}}]{
                \texttt{\tiny{[GM$\langle$P2]}}
                \texttt{ΑΠΟΦΑΣΗ: διαφορετικά πλέγματα} \\
            }
        }
    }
     \\ \\

    \theutterance \stepcounter{utterance}  
    & & & \multicolumn{2}{p{0.3\linewidth}}{
        \cellcolor[rgb]{0.9,0.9,0.9}{
            \makecell[{{p{\linewidth}}}]{
                \texttt{\tiny{[GM$|$GM]}}
                \texttt{continue} \\
            }
        }
    }
    & & \\ \\

    \theutterance \stepcounter{utterance}  
    & & & \multicolumn{2}{p{0.3\linewidth}}{
        \cellcolor[rgb]{0.9,0.9,0.9}{
            \makecell[{{p{\linewidth}}}]{
                \texttt{\tiny{[GM$|$GM]}}
                \texttt{success} \\
            }
        }
    }
    & & \\ \\

    \theutterance \stepcounter{utterance}  
    & & \multicolumn{4}{p{0.6\linewidth}}{
        \cellcolor[rgb]{0.9,0.9,0.9}{
            \makecell[{{p{\linewidth}}}]{
                \texttt{\tiny{[P1$\langle$GM]}}
                \texttt{ΑΠΑΝΤΗΣΗ: Όχι, δεν υπάρχει 'X' στη θέση (2,1).} \\
\texttt{Τώρα πάρε μια απόφαση. Τι πιστεύεις: είναι το δικό σου πλέγμα και το άλλο πλέγμα που περιγράφεται το ίδιο; Γράψε "ΑΠΟΦΑΣΗ: ίδιο πλέγμα" αν πιστεύεις ότι είναι το ίδιο πλέγμα ή "ΑΠΟΦΑΣΗ: διαφορετικά πλέγματα" αν πιστεύεις ότι είναι διαφορετικά πλέγματα. Μην προσθέσεις τίποτα άλλο.} \\
            }
        }
    }
    & \\ \\

    \theutterance \stepcounter{utterance}  
    & \multicolumn{4}{p{0.6\linewidth}}{
        \cellcolor[rgb]{0.9,0.9,0.9}{
            \makecell[{{p{\linewidth}}}]{
                \texttt{\tiny{[P1$\rangle$GM]}}
                \texttt{ΑΠΟΦΑΣΗ: διαφορετικά πλέγματα} \\
            }
        }
    }
    & & \\ \\

    \theutterance \stepcounter{utterance}  
    & & & \multicolumn{2}{p{0.3\linewidth}}{
        \cellcolor[rgb]{0.9,0.9,0.9}{
            \makecell[{{p{\linewidth}}}]{
                \texttt{\tiny{[GM$|$GM]}}
                \texttt{continue} \\
            }
        }
    }
    & & \\ \\

    \theutterance \stepcounter{utterance}  
    & & & \multicolumn{2}{p{0.3\linewidth}}{
        \cellcolor[rgb]{0.9,0.9,0.9}{
            \makecell[{{p{\linewidth}}}]{
                \texttt{\tiny{[GM$|$GM]}}
                \texttt{success} \\
            }
        }
    }
    & & \\ \\

\end{supertabular}
}

\end{document}
